% Supplementary Information for Nature submission
\documentclass[11pt]{article}

\usepackage[utf8]{inputenc}
\usepackage[T1]{fontenc}
\usepackage{mathptmx}
\usepackage{amsmath,amssymb,amsthm}
\usepackage{graphicx}
\usepackage[margin=2.5cm]{geometry}
\usepackage{setspace}
\usepackage{hyperref}
\usepackage{booktabs}
\usepackage{caption}
\usepackage{float}
\usepackage{siunitx}

\doublespacing

\newtheorem{theorem}{Theorem}
\newtheorem{proposition}{Proposition}
\newtheorem{lemma}{Lemma}

\renewcommand{\thefigure}{S\arabic{figure}}
\renewcommand{\thetable}{S\arabic{table}}
\renewcommand{\theequation}{S\arabic{equation}}
\renewcommand{\thesection}{S\arabic{section}}

\newcommand{\Vb}{V_{\mathrm{b}}}
\newcommand{\Vd}{V_{\mathrm{d}}}
\newcommand{\DV}{\Delta V}

\title{\textbf{Supplementary Information}\\[0.5em]
A unified control equation for cell size homeostasis across domains of life}

\author{}
\date{}

\begin{document}

\maketitle
\tableofcontents
\clearpage

\section{First-principles derivation of the Hill-type hazard function}

\subsection{Molecular setup and concentration-dilution sensor}

Consider a cell of volume $V(t)$ growing exponentially at rate $\mu$.
The cell contains a division-inhibiting molecule $I$ (e.g., Whi5 in
\textit{S.~cerevisiae}, Rb in mammalian cells) whose total amount is
synthesized once per cell cycle: $I_{\mathrm{tot}} = \text{const}$.
As the cell grows, the inhibitor concentration decreases:
\begin{equation}
  [I](V) = \frac{I_{\mathrm{tot}}}{V}
\end{equation}

Division is promoted by a cyclin-CDK activator complex $A$ at constant
concentration $[A]_0$.  The activator targets respond cooperatively to
the ratio of activator to inhibitor signals through $n$ cooperative
binding sites:
\begin{equation}
  f_{\mathrm{active}}(V) = \frac{[A]_0^n}{[A]_0^n + [I]^n}
    = \frac{V^n}{V^n + K^n}
\end{equation}
where $K \equiv I_{\mathrm{tot}} / [A]_0$ is the volume at which
inhibitor and activator signals balance.

The instantaneous division rate (hazard) is $f_{\mathrm{active}}$
times a maximal division attempt rate $\mu r$:
\begin{equation}
  h(V) = \mu \, r \, \frac{V^n}{V^n + K^n}
\end{equation}

\textbf{Physical interpretation of parameters}:
$K$ = volume at half-maximal activation, set by inhibitor dosage
($K = I_{\mathrm{tot}}/[A]_0$), scales with ploidy and gene dosage;
$n$ = cooperativity of the molecular switch (number of phosphorylation
sites or binding events);
$r$ = maximal division rate in units of $\mu^{-1}$.

\subsection{Closed-form survival function}

The survival function admits a closed form via the substitution
$u = V'^n + K^n$:
\begin{equation}
  S(V \mid \Vb) = \exp\!\left(-r \int_{\Vb}^{V}
    \frac{V'^{n-1}}{V'^n + K^n}\, dV'\right)
  = \left(\frac{\Vb^n + K^n}{V^n + K^n}\right)^{r/n}
\end{equation}

The division volume density is:
\begin{equation}
  p(V_d \mid \Vb) = \frac{r\, V_d^{n-1}}{V_d^n + K^n}
    \left(\frac{\Vb^n + K^n}{V_d^n + K^n}\right)^{r/n}
\end{equation}

\subsection{Justification of the $\mu$ prefactor}

The factor $\mu$ ensures dimensional consistency and captures the
empirical growth-rate dependence of division.  In \textit{E.~coli},
chromosome replication initiation is triggered by DnaA accumulation,
whose production rate is proportional to growth
rate\cite{si2019}.  The division machinery itself (FtsZ ring in
bacteria, cytokinetic ring in eukaryotes) is growth-rate dependent.

\section{Existence and uniqueness of the stationary size distribution}

\begin{proposition}[Existence of stationary distribution]
For any $n > 0$, $K > 0$, $r > 0$, the iterated map
$\Vb^{(g+1)} = \Vd^{(g)}/2$, where $\Vd^{(g)}$ is drawn from
$p(\Vd \mid \Vb^{(g)})$, admits a stationary distribution $\pi(\Vb)$
on $(0, \infty)$.
\end{proposition}

\begin{proof}[Proof sketch]
The transition kernel $T(\Vb' \mid \Vb) = 2\, p(2\Vb' \mid \Vb)$
maps the space of probability distributions on $(0,\infty)$ to itself.
We verify:

\textbf{(i) Tightness.}  For any $\Vb$, the division distribution
$p(\Vd \mid \Vb)$ has support on $[\Vb, \infty)$ with exponentially
decaying tail (since $S(\Vd) \leq (\Vb/\Vd)^{r/2}$ for all
parameter values).  Thus the sequence of iterates $\{\Vb^{(g)}\}$ is
tight.

\textbf{(ii) Irreducibility.}  For any two intervals $(a,b)$ and
$(c,d)$ with $0 < a < b$ and $0 < c < d$, there exists $g^*$ such
that $P(\Vb^{(g^*)} \in (c,d) \mid \Vb^{(0)} \in (a,b)) > 0$.  This
follows because the division density $p(\Vd \mid \Vb) > 0$ for all
$\Vd > \Vb$.

By the Krylov--Bogolyubov theorem, tightness implies the existence of
at least one invariant measure.
\end{proof}

\begin{proposition}[Uniqueness and geometric ergodicity]
The stationary distribution $\pi$ is unique, and the iterates
converge geometrically: there exist constants $C > 0$ and
$\rho \in (0,1)$ such that
$d_W(\pi^{(g)}, \pi) \leq C \rho^g$ for any initial distribution
$\pi^{(0)}$, where $d_W$ is the Wasserstein-1 distance.
\end{proposition}

\begin{proof}[Proof sketch]
We establish a contraction in the log-volume coordinate
$u = \ln V$.  The transition kernel in $u$-space has the form
$u' = \ln(\Vd/2)$ where $\Vd$ is drawn from a distribution that
concentrates near $\Vb \cdot 2^{1/r}$ (the median).  For the
adder regime ($n=1$, $K \gg \Vb$), the map $\langle u' \rangle
\approx \ln(K \ln 2 / (2r))$ is independent of $u$, giving strong
contraction.  For the sizer regime ($n \to \infty$), $\langle u'
\rangle \approx \ln(K/2)$, again independent of $u$.  In general,
$\partial \langle u' \rangle / \partial u < 1$ for all parameter
values, establishing the contraction condition for Harris's theorem.
\end{proof}

\section{Detailed limiting-case derivations}

\subsection{Sizer limit ($n \to \infty$)}

For $V < K$: $V^n/K^n = (V/K)^n \to 0$, so $h(V) \to 0$.

For $V > K$: $V^n/(V^n + K^n) = 1/(1 + (K/V)^n) \to 1$, so
$h(V) \to \mu r$.

The survival function for $V \geq K$:
\begin{equation}
  S(V) = \exp\left(-r \int_K^V \frac{dV'}{V'}\right)
       = \left(\frac{K}{V}\right)^r
\end{equation}

As $r \to \infty$, $S(V) \to 0$ for any $V > K$, and the division
distribution becomes $p(\Vd) \to \delta(\Vd - K)$.  Division size is
$K$ regardless of $\Vb$: \textbf{pure sizer}.

\textbf{Error estimate:} For large but finite $n$, the transition
width from $h \approx 0$ to $h \approx \mu r$ scales as
$\delta V / K \sim 1/n$ (from the derivative of the Hill function at
$V = K$).  The CV of division size due to this finite sharpness is
$\mathrm{CV} \sim 1/(nr)$.

\subsection{Timer limit ($n \to 0^+$)}

As $n \to 0^+$, for any $V > 0$:
$V^n = e^{n \ln V} \to 1$ and $K^n \to 1$.
Therefore:
\begin{equation}
  h(V) \to \mu r \cdot \frac{1}{1+1} = \frac{\mu r}{2}
\end{equation}

This is a constant hazard (independent of $V$).  Converting to time
using $V = \Vb e^{\mu t}$:
\begin{equation}
  S(t) = \exp\left(-\frac{r}{2} \mu t\right)
\end{equation}

This is an exponential distribution with rate $\lambda = \mu r / 2$.
The mean interdivision time $\langle T \rangle = 2/(\mu r)$ is
\textbf{independent of birth size}: a stochastic timer.

\textbf{Important caveat:} This is a \emph{memoryless} (exponential)
timer, not a deterministic timer with fixed $T$.  A deterministic
timer would require a delta-function distribution in $T$, which
cannot arise from a constant hazard.  The exponential timer has
$\mathrm{CV}(T) = 1$, whereas a deterministic timer has
$\mathrm{CV}(T) = 0$.  Biological timers (e.g., light-cycle control
in \textit{Chlamydomonas}\cite{donnan1983}) may involve additional
clock mechanisms not captured by the hazard framework.

\textbf{Error estimate:} For small but finite $n$, the hazard acquires
a weak size dependence: $h(V) \approx (\mu r/2)(1 + n \ln(V/K)/2 +
O(n^2))$.  The induced correlation between $T$ and $\Vb$ scales as
$|\mathrm{Corr}(T, \Vb)| \sim n/2$.

\subsection{Adder limit ($n = 1$, $K \gg \langle\Vb\rangle$)}

With $n = 1$, the survival function becomes:
\begin{equation}
  S(V \mid \Vb) = \left(\frac{\Vb + K}{V + K}\right)^r
\end{equation}

The division density is a shifted Pareto-like distribution:
\begin{equation}
  p(\Vd \mid \Vb) = \frac{r}{\Vd + K}
    \left(\frac{\Vb + K}{\Vd + K}\right)^r, \qquad \Vd \geq \Vb
\end{equation}

This distribution is \emph{monotonically decreasing} on
$[\Vb, \infty)$: its logarithmic derivative is
$-(r+1)/(\Vd+K) < 0$ everywhere.  The mode is at the boundary
$\Vd = \Vb$, not at an interior point.

The \textbf{mean} division volume is computed exactly:
\begin{equation}
  \langle \Vd \rangle = (\Vb + K) \cdot \frac{r}{r-1} - K
  = \frac{\Vb + K}{r-1}
  \quad (\text{for } r > 1)
\end{equation}
giving $\langle \DV \rangle = (\Vb + K)/(r-1)$.

The \textbf{median} satisfies $S(\Vd^{\mathrm{med}}) = 1/2$:
\begin{equation}
  \Vd^{\mathrm{med}} = (\Vb + K) \cdot 2^{1/r} - K
\end{equation}

For $K \gg \Vb$, both converge to constants independent of $\Vb$:
\begin{equation}
  \langle \DV \rangle \approx \frac{K}{r-1}
    \left[1 + \frac{\Vb}{K} + O\!\left(\frac{\Vb^2}{K^2}\right)\right]
\end{equation}
which is the \textbf{pure adder}: added volume independent of birth
size to leading order.

\textbf{Error estimate:} The relative deviation from perfect adder is
$|\DV(\Vb) - \DV(0)|/\DV(0) = \Vb/K$.  For \textit{E.~coli},
$\langle\Vb\rangle \approx 1.5\,\mu\mathrm{m}^3$ and
$K \approx 15\,\mu\mathrm{m}^3$ (from fits), giving $\sim 10\%$
deviation---consistent with the experimentally observed weak residual
slope in the $(\Vb, \DV)$ plot.

\section{Transient overshoot derivation (Prediction 1)}

After a nutrient upshift at $t = 0$, $\mu$ jumps from $\mu_1$ to
$\mu_2 > \mu_1$, while $K$ adjusts on a slower timescale $\tau_K$:
\begin{equation}
  K(t) = K_2 - (K_2 - K_1)e^{-t/\tau_K}
\end{equation}

Since the hazard scales with $\mu$, the growth and division rates
both increase immediately.  However, if $K_2 > K_1$ (faster growth
requires larger cells, as observed empirically), $K$ lags behind.

\begin{proposition}[Overshoot condition]
A transient overshoot in mean cell size occurs when:
\begin{equation}
  \tau_K > \frac{1}{\mu_2}\cdot\frac{n}{r}\cdot
    \frac{K_2}{K_2 - K_1}
\end{equation}
The overshoot magnitude is:
\begin{equation}
  \frac{\langle V \rangle_{\max} - \langle V \rangle_{\mathrm{ss}}}
       {\langle V \rangle_{\mathrm{ss}}}
  \approx \frac{K_2 - K_1}{K_2}\cdot\left(1 - e^{-\mu_2\tau_K}\right)
\end{equation}
\end{proposition}

The overshoot arises because the effective division size temporarily
decreases (old $K$) while cells grow faster (new $\mu$), creating a
mismatch that resolves non-monotonically as $K(t) \to K_2$.

\section{Phase-specific control combination rule (Prediction 3)}

% TODO: Analytical derivation of n_eff from phase-specific n values.
% To be completed with input from cell-size-physicist.

For a cell cycle divided into two phases (e.g., G1 and S/G2/M), each
governed by its own hazard function $h_1(V)$ and $h_2(V)$ with Hill
coefficients $n_1$ and $n_2$ respectively, the whole-cycle effective
behaviour depends on the \emph{sequential} application of the two
hazard functions.

In budding yeast, the G1 phase operates as a sizer ($n_1 \gg 1$,
driven by Whi5 dilution) and the post-Start phase as a timer
($n_2 \approx 0$).  The composite produces approximate adder
behaviour at the whole-cycle level, as observed by
Chandler-Brown \textit{et al.}\cite{chandler2017}.

\section{Population-level dynamics and moment equations}

\subsection{Structured population equation (Fokker--Planck)}

Let $\phi(V, t)$ be the volume density of cells at time $t$.  The
population evolves according to:
\begin{equation}
  \frac{\partial \phi}{\partial t} + \frac{\partial}{\partial V}(\mu V \phi)
  = -h(V)\phi + 4\mu r \frac{(2V)^n}{(2V)^n + K^n}\phi(2V, t)
\end{equation}
where the first term on the right is loss from division at rate
$h(V)$, and the second is gain from mothers of size $2V$ dividing
symmetrically.  At steady state, this reduces to the renewal equation
in Section~S2, confirming internal consistency.

This population-level PDE is the deterministic ``master equation''
of the system.  The hazard function $h(V)$ encodes all control
information: the framework is a stochastic process at the single-cell
level whose population statistics are governed by a structured
population equation.

\subsection{Moment equations and the two-parameter reduction}

Let $\rho_g(\Vb)$ be the birth-size density in generation $g$.  The
map $\Vb^{(g+1)} = \Vd^{(g)}/2$ gives:
\begin{equation}
  \rho_{g+1}(\Vb') = 2 \int_0^\infty p(2\Vb' \mid \Vb)\,
    \rho_g(\Vb)\, d\Vb
\end{equation}

At steady state, the condition $\langle \Vd \rangle = 2\langle \Vb \rangle$
(cells must double on average for population stability) constrains $r$
as a function of $n$ and $K/\langle\Vb\rangle$.  For $n = 1$:
$r = 1 + K/\langle\Vb\rangle \cdot (1/\ln 2)$.

This explains why the model has effectively \textbf{two free control
parameters} $(n, K)$ despite having three: $r$ is determined by the
steady-state constraint.  The ``two-parameter family'' refers to the
$(n, K)$ plane that defines the phase diagram of control strategies,
where $n$ sets the \emph{type} of control (cooperativity) and $K$
sets the \emph{target size} (inhibitor dosage).  The third parameter
$r$ controls \emph{precision} (fidelity of control) and is fixed by
the doubling requirement.

\section{Supplementary Figures}

% Placeholder for supplementary figures.
% To be populated after BioPlotAnalyst produces figures.

% \begin{figure}[H]
% \centering
% \includegraphics[width=0.8\textwidth]{../figures/supp_fig1.pdf}
% \caption{\textbf{Supplementary Figure S1 $|$ Convergence of the
% iterated map.} Starting from three different initial distributions
% (delta, uniform, log-normal), the birth-size distribution converges
% to the same stationary distribution within $\sim 30$ generations
% for all parameter regimes tested.}
% \end{figure}

\section{Supplementary Tables}

% Placeholder for parameter tables.
% To be populated after CodeAgent005 produces fitting results.

% \begin{table}[H]
% \centering
% \caption{\textbf{Supplementary Table S1 $|$ Complete fitted
% parameters.} Maximum-likelihood estimates of $(n, K, r)$ for all
% species, conditions, and replicates, with 95\% confidence intervals.}
% \begin{tabular}{lccccc}
% \toprule
% Species & Condition & $n$ & $K$ & $r$ & AIC \\
% \midrule
% & & & & & \\
% \bottomrule
% \end{tabular}
% \end{table}

\end{document}
